\subsection{An explicit order-128 obstruction for an ordinary class-four ghost}

\begin{statement}{A.207}{Proved}
{A prescribed ordinary residue, an actual global class-four lift, and its finite Selmer obstruction}
There is an infinite-support ordinary completed unit-equation pair
\((A,B)\) with
\[
 (A,B)\bmod8=([3],[-2]),\qquad
                  t_{1129}\equiv828\pmod {1129},
\tag{A.207.1}
\]
where \((t_p,1-t_p)\) denotes its ordinary localization.
The pair has an actual global section in
\(E_4=\pi_1(U)/\Gamma_5\), and therefore also one in
\(E_3\).  It has no everywhere locally point-or-cusp-matching
section in \(E_4\).

The obstruction uses the actual quotient \(Q_1\) of A.193,
of geometric order \(128\) and nilpotency class four.
More precisely, no everywhere locally point-or-cusp-matching
section of that arithmetic quotient in the combined eighth-power
class \([-24]\) can have an ordinary local image at a point
\[
                         t\in828+1129\Z_{1129}.
\tag{A.207.2}
\]
Every one of the eight rational local cover sheets is excluded
by its explicit two-Kummer class below.  The construction gives
continuous global class-three and class-four lifts; it does not
assert all-place local matching for its class-three lift.
\end{statement}

\begin{proof}
\emph{A literal first-stage specialization of the existing recursion.}
Use A.62 with \(C=24\), \(J=6\), so that
\[
 \begin{aligned}
 n_j&=(j+6)!,&D_j&=144n_j^3,&
 E_j(q)&=24n_j(q^2-1)(q^3-1),\\
 A&=3\prod_{j\ge0}\frac{q_j^{E_j(q_j)}}{s_j^{E_j(s_j)}},&
 B&=-2\prod_{j\ge0}\frac{r_j^{E_j(r_j)}}{s_j^{E_j(s_j)}}.
 \end{aligned}
\tag{A.207.3}
\]
Only one local choice in the first triangular CRT step is
changed.  Put \(p=1129\).  The exact initial data are
\[
 \begin{gathered}
 p-1=24\cdot47,\quad n_0=720,\quad D_0=53747712000,
 \quad\gcd(D_0,p-1)=24,\\
 59^{D_0}\equiv276,\qquad28^{D_0}\equiv978,\qquad
                    3\cdot276-2\cdot978\equiv1\pmod p.
 \end{gathered}
\tag{A.207.4}
\]
The derivative of \(Z\mapsto Z^{D_0}\) at one is a
\(p\)-unit.  Hence there is a unique
\(Z_0\in1+p\Z_p\) satisfying
\[
             Z_0^{D_0}=3\cdot59^{D_0}-2\cdot28^{D_0}.
\tag{A.207.5}
\]
This gives an actual unit-coordinate \(D_0\)-point on the
initial homogeneous curve, with the displayed power residues.

When the triangular procedure reaches \(p\), all residues
at prime divisors of \(p-1\) have already been fixed.  They
determine the three actual exponent classes \(e_i\bmod(p-1)\)
in (A.62.11).  That identity gives
\(\gcd(e_i,p-1)=24\), independently for all three variables.
Their common residue power image therefore has order \(47\).
Each of \(276,978,1\) lies in it, by (A.207.4), and each
has exactly twenty-four preimages under each exponent map.
At most four nonzero residues satisfy \(u^2=1\) or \(u^3=1\).
Thus each target has an admissible preimage satisfying the
literal exclusions (A.62.10).

On its chosen residue ball the actual map (A.62.13) is a
bijection onto that Teichmuller value times \(1+p\Z_p\).
Indeed the derivative coefficient
\(24n_0(\omega^2-1)(\omega^3-1)\) in (A.62.14) is a
\(p\)-unit.  Choose its three input values so that its
outputs are exactly
\[
                  59^{D_0},\quad28^{D_0},\quad Z_0^{D_0}.
\tag{A.207.6}
\]
Their homogeneous equation is precisely (A.207.5).
Approximate those inputs to the original precision \(p^{L_p}\)
of A.62.  The exact local equation and its continuity give
the same next-stage error bound as before.  All three inputs
still satisfy the exclusions that the later triangular gcd
argument uses.  Thus the change affects no assertion about
the exponent classes at subsequent processed primes.

The three complete-splitting fields, their intersections with
the cyclotomic CRT field, and the unit classes equal to one
modulo \(n_0\) are unchanged.  The actual Chebotarev choice
of the new primes in A.59.4--6 therefore still applies.
The new coefficients satisfy
\[
       (a_1,b_1,c_1)\equiv(828,302,1)\pmod p,
                  \qquad a_1/c_1\equiv828\pmod p.
\tag{A.207.7}
\]
The error precision exceeds the next Hensel threshold, so
the next all-place homogeneous-unit-point invariant remains
valid at \(p\), as at all other old and new primes.

For all later stages \(p\le B(D_j)\).  If \(p\nmid n_j\),
use the original small-prime construction in A.62: the three
power values have one common nonzero residue.  Their slight
principal-unit adjustments do not change that residue.
Multiplying the three coefficients by these values leaves
\(a_j/c_j\) unchanged modulo \(p\).
Once \(p\mid n_j\), all three prescribed input balls are
one modulo \(p\), and their power values are one modulo
\(p\) as well.  This again preserves the ratio.
Every new coefficient prime exceeds the processed-prime bound,
so none equals \(p\).  The coefficients at \(p\) remain units.
The ordinary limits proved in A.62 consequently satisfy
\(t_p\equiv828\pmod p\).

All other recursion inputs and proofs of A.62 are retained:
the local limits are ordinary at every finite place; the
supports are disjoint and infinite; and the valuation pairs
at two and three remain \((0,1)\) and \((1,0)\).
Every update exponent is divisible by twenty-four, giving
the two individual mod-eight classes in (A.207.1).
The exponents in (A.207.3) have exactly the form (A.61.1),
with the factorial coefficients tending to zero in
\(\widehat\Z\).  Applying A.61 therefore supplies an actual
global class-four section above this pair.

This specifies an effective prime recursion.  For example,
choose the least admissible residue preimages, the unique
Hensel lifts on their prescribed balls, and then the least
successive primes in the established splitting/CRT classes.
The finite data (A.207.4) and all preimage checks have been
executed.  No full, very large first-stage prime search is
claimed to have been executed.

\emph{The exact finite geometric quotient and its Selmer map.}
Use the actual cover and proper compactification
\[
 \begin{aligned}
 W&:\quad z^8=-\frac{t(1-t)^3}{24},\\
 C&:\quad y^2=x^5-54x,\\
 t&=\frac{x^4}{x^4-54},\quad z=\frac{3x}{y},\qquad
 x=-\frac{6z^2}{1-t},\quad y=-\frac{18z}{1-t}.
 \end{aligned}
\tag{A.207.8}
\]
Both displayed function-field maps and their inverse
compositions were proved in (A.192.12).  Let
\(H\subset\Pi\) be the index-eight geometric subgroup,
and let \(K\) be the kernel of the actual composition
\[
       H\longrightarrow\pi_1(C_{\overline\Q})^{\rm ab}
                    \longrightarrow J[2],
       \qquad J=\operatorname{Jac}(C).
\tag{A.207.9}
\]
The rational origin \(\infty\) fixes the Abel--Jacobi map.
Deck transformations and arithmetic automorphisms preserve
compactification and the mod-two kernel, so \(K\) is
normal in the arithmetic group.

For clarity, the literal group calculation of A.193 at
\(n=1\) is small.  In the Nielsen generators
\(v=x, u=yx^{-3}\), the deck character is zero on
\(u\) and one on \(v\).  Its Schreier generators are
\(u_j=v^juv^{-j}\), \(0\le j<8\), and \(z=v^8\).
In proper mod-two homology the boundary relations give
\(z=0\) and \(u_{j+4}=u_j\); the remaining sum relation
is then redundant.  Thus, writing \(\delta=v-1\),
\[
 \begin{aligned}
 Q_1=\Pi/K&\simeq
       \bigl(\F_2[\delta]/\delta^4\bigr)\rtimes C_8,\\
 x&\longmapsto v,\qquad y\longmapsto u_0v^3,\\
 |Q_1|&=128,\qquad
 \Gamma_j(Q_1)=\delta^{j-1}\F_2[\delta]/\delta^4
                       \quad(2\le j\le5).
 \end{aligned}
\tag{A.207.10}
\]
Here \(v\) acts by \(1+\delta\), and the last group is
zero at \(j=5\).  The term at \(j=4\) contains the
nonzero word \([[[u_0,v],v],v]=\delta^3u_0\).
This proves the class-four assertion and the actual map
\(E_4\to\pi_1(U)/K\), since \(\Gamma_5\Pi\subset K\).
The full finite group and all lower-central orders were
also checked directly.

Consider an everywhere locally matching section of this
arithmetic quotient in the stated combined eighth-power
class.  Pulling back the torsor in (A.207.8) gives the
trivial Kummer class: its radicand is
\(-t(1-t)^3/24\), and at the global classes of (A.207.1)
it becomes \(-3(-2)^3/24=1\) modulo eighth powers.
Select its rational fixed sheet.  The resulting lift to
the arithmetic cover subgroup, followed by compactification
and Abel--Jacobi modulo two, gives an actual class
\(\xi\in H^1(G_\Q,J[2])\).

At every place, align the given quotient section with its
actual full local point or cusp image and transport the
selected finite sheet by that geometric conjugator.
For an ordinary point the fixed sheet is a rational point
of the finite etale fibre in \(W\).  For a cusp the same
fixed-sheet decomposition-group map compactifies to the
rational boundary point of \(C\).  The real comparison
uses the same rational sheet and proper real centre.
These are the finite-cover maps of A.178 and A.193; they
show that \(\xi\) belongs to \(\operatorname{Sel}^2(J/\Q)\).
No choice of a local sheet has been discarded.

The exact global calculation (A.186.1), (A.186.18)--21
identifies this Selmer group with the four classes generated
by
\[
 D=[(3,9)-\infty],\quad T=[(0,0)-\infty],\qquad
 \kappa_2(D)=3-\theta,\quad\kappa_2(T)=-\theta,
                         \quad\theta^4=54.
\tag{A.207.11}
\]
These are the actual odd-degree Kummer coordinates, with
the rational linear factor restored by
\(\eta\mapsto(N\eta,\eta)\) as in A.186.
There is no further diagonal-square identification.

\emph{All eight local sheets and the evaluated obstruction.}
At \(p=1129\), the four simple quartic roots are
\[
                 (16,430,699,1113)\pmod p.
\tag{A.207.12}
\]
They lift uniquely to the four roots in \(\Q_p\).
There is an explicit geometric basis behind the four
local coordinates.  Write \(T_i=[(\theta_i,0)-\infty]\)
for the four nonrational branch classes and retain \(T\)
for the class at zero.  The branch Weil pairing of A.186
gives
\[
 e_i=T+T_i,\qquad e_2(e_i,T_j)=(-1)^{\delta_{ij}}.
\tag{A.207.13}
\]
Thus the squareclass coordinates \(x-\theta_j\) are
literally the coordinates of the Kummer cocycle in the
dual basis \(e_i\).  Under the deck map
\((x,y)\mapsto(ix,\zeta_8y)\), this basis cycles with
the roots.  At \(1129\), \(i=168\) and its cyclic root
order is \((16,430,1113,699)\); equivalently it acts
as the cycle \((1\ 2\ 4\ 3)\) in (A.207.12).
This specifies the literal kernel action in (A.207.10)
without a choice of an abstract four-dimensional vector
space.  At this split prime, unit quadratic characters
are unramified, and their bits are exactly their
Frobenius cocycle values in this basis.

For a unit in the corresponding split quartic algebra,
record zero for a square residue and one for a nonsquare
residue, in this order of embeddings.  The literal
localization of (A.207.11) is
\[
                 \kappa_2(T)=0000,\qquad
                 \kappa_2(D)=1111.
\tag{A.207.14}
\]
Thus the complete localized global two-Selmer image is
\(\{0000,1111\}\).  The entries here have valuation zero;
the recorded unit-square classes are their full local
squareclass data.

For \(t\equiv828\), the right side of the eighth-root
equation in (A.207.8) is \(1115\) modulo \(p\).
Its eight simple roots and the associated compact points are
\[
\begin{array}{c|r|r|c}
 z&x&y&\kappa_2([(x,y)-\infty])\\\hline
 35&28&454&0101\\
 44&941&603&1100\\
 235&1101&629&1010\\
 511&188&306&0011\\
 618&188&823&0011\\
 894&1101&500&1010\\
 1085&941&526&1100\\
 1094&28&675&0101
\end{array}
\tag{A.207.15}
\]
The last column is obtained by applying the four unit
quadratic characters to \(x-\theta\), exactly the point
Kummer map of A.186.  All its entries are units, and no
displayed vector belongs to (A.207.14).
The derivatives of the eighth-root and quartic polynomials
are units at these roots.  Hensel's lemma therefore gives
all eight rational sheets over every \(t\) in (A.207.2),
and their unit-square vectors are unchanged throughout
that residue disc.  This proves the stated obstruction
for the arithmetic quotient.

Finally an all-place locally matching \(E_4\)-section
above the pair constructed in (A.207.3) would map to this
quotient and have a trivial pulled-back eighth-root class.
At \(1129\), full local multiplicative Kummer injection
forces its ordinary point centre to be the specified
\(t_{1129}\).  No cusp can have that abelian pair:
the three regular cusp coordinates would require,
respectively, \(1-t=1\), \(t=1\), or
\((1-t)/t=-1\).  Thus the local image lies in (A.207.2),
contrary to (A.207.15).

The contradiction uses the global Selmer image supplied by
all local comparisons and then evaluates it at \(1129\).
It does not claim that every unrestricted global quotient
section is already excluded at that place alone.
In particular every one of the \(240^3\) global
class-four corrections above the constructed lower
class-three representative fails the all-place matching
test of A.205, while no all-place matching of that lower
representative has been assumed or proved.
\end{proof}

The exact finite certificate
\path{code/point3_class_four_1129_certificate.py} and its
retained output and data verify primality, the initial
homogeneous power residues, all 138 exponent/target
preimage cases, the later residue-ratio invariance,
the eight actual cover sheets and their Kummer vectors,
and the order and lower central series of the literal
finite quotient.  The infinite recursion and its
continuous section are supplied by the proved modifications
above, not by an unperformed large numerical prime search.
